\section{Models and conventions}
\label{sec:models}

This section fixes the input encodings, the promises and certificates that
may accompany an input, the output contracts, and the computational models.
Later sections use these definitions without restating them. The section
also proves a few elementary facts that are used throughout.

\subsection{Notation}
\label{sec:models-notation}

For $x\in\R^n$, $\norm{x}$ is the Euclidean norm, and $\norm{x}_1$ and
$\norm{x}_\infty$ are the $\ell_1$ and maximum norms. For a set
$S\subseteq\R^n$, $\dist(x,S)=\inf_{s\in S}\norm{x-s}$. The space of real
symmetric $m\times m$ matrices is $\mathrm{Sym}_m$; $A\succeq0$ and $A\succ0$
mean positive semidefinite and positive definite, and $\lambda_{\min}(A)$ is
the least eigenvalue. The Kronecker product $x\otimes v\in\R^{n^2}$ has
coordinates $x_iv_j$, ordered lexicographically by $(i,j)$.

The following symbols are reserved. The total binary input length is $L$
(Definition~\ref{def:models-length}); the number of continuous variables is
$n$; $k$ is a structural parameter defined in each statement; $q$ is a
number of requested bits of precision, and $\varepsilon=2^{-q}$. The objective
is $f$, the feasible set $P$, the optimal value $f_*=\inf_Pf$, and the set of
optimizers $S$. The selected optimizer is $p$; it is the unique optimizer
when $f$ is strictly convex on a convex feasible set, and otherwise, for
example in mixed-integer problems with several optimal points, the selection
rule is stated. An
explicit polynomial whose sign or value at $p$ is requested is an
\emph{observable} and is denoted by $h$. Lower curvature bounds are denoted
by $\mu$ and upper curvature or Lipschitz bounds by $\Lambda$. The degree of
the input polynomials is $D$. We write $\poly(\cdot)$ for a polynomial whose
exponent is an absolute constant, $C$ for an absolute constant, and $F(k)$,
$a(t)$ for computable functions of the parameters shown. For a point $p$ with
algebraic coordinates, $\Q(p)=\Q(p_1,\ldots,p_n)$ is its coordinate field.

\subsection{Input encodings}
\label{sec:models-input}

\begin{definition}[Rational data]
\label{def:models-rational}
A rational number $a/b$ with $b\ge1$ and $\gcd(a,b)=1$ is encoded by $a$ and
$b$ in binary. Its bit length is
$\lceil\log_2(\abs a+1)\rceil+\lceil\log_2(b+1)\rceil+1$, and the bit length
of a vector or matrix is the sum of the bit lengths of its entries.
\end{definition}

\begin{definition}[Explicit polynomials and degree]
\label{def:models-polynomial}
An \emph{explicit} (or explicitly sparse) polynomial
$f=\sum_\alpha f_\alpha x^\alpha\in\Q[x_1,\ldots,x_n]$ is given by the list
of its nonzero coefficients $f_\alpha$ together with their exponent vectors
$\alpha\in\N^n$. Its degree is $\deg f=\max\{\abs\alpha:f_\alpha\ne0\}$. When
the degree is not fixed, we use the numerical degree $D=\max\{2,\deg f\}$, or
the maximum over all input polynomials. A bound of the form $\poly(L,D)$ is
polynomial in the input length when the degree is at most polynomial in
$L$, for example when exponents are written in unary; we then speak of
\emph{unary degree}.
\end{definition}

The \emph{interaction graph} of a polynomial has one vertex for each
variable and an edge between two variables when a nonzero monomial contains
both. For a polynomial, this is equivalent to the corresponding off-diagonal
entry of its Hessian being a nonzero polynomial. The graph records which
variables interact; bounds on its treewidth or a supplied forest are stated
where used.

\begin{definition}[Polyhedra]
\label{def:models-polyhedra}
A rational polyhedron is $P=\{x\in\R^n:Bx\le b\}$ with $B\in\Q^{m\times n}$
and $b\in\Q^m$; an equality is written as two inequalities. A box is
$\{x:l\le x\le u\}$ with $l_i\in\Q\cup\{-\infty\}$ and
$u_i\in\Q\cup\{+\infty\}$. A network-flow polyhedron is
$\{x\in\R^E:Ax=b,\ l\le x\le u\}$, where $A$ is the node-arc incidence
matrix of a directed graph with arc set $E$.
\end{definition}

\begin{definition}[Input length]
\label{def:models-length}
The input length $L$ is the total bit length of all input data: the
polynomials (objective, constraints, observables), the polyhedron,
thresholds, supplied numbers such as curvature bounds, slack gaps, or noise
widths, and supplied certificates. The number of variables is part of the
input, so $n\le L$, and we assume $L\ge2$.
\end{definition}

Quantities computed from the input in polynomial time have bit length
$\poly(L)$, so bounds stated in terms of $L$ remain valid when such
quantities are added to the input.

\begin{remark}[Encodings not covered]
\label{rem:models-encodings}
Three other encodings are excluded unless a statement says otherwise. First,
a polynomial given by an arithmetic circuit can have exponentially many
monomials and exponentially large degree. Second, exponents written in binary
allow degree exponential in $L$; our bounds that are polynomial in $D$ are
then not polynomial in $L$. Third, expansion is not free under substitution:
$(x_1+\cdots+x_n)^D$ has $\binom{n+D-1}{D}$ monomials, so a rational affine
substitution $x=x_0+Vu$ can make a short polynomial of growing degree long.
Algorithms for variable degree therefore evaluate such compositions instead of
expanding them. At fixed degree, expansion changes the length only
polynomially.
\end{remark}

\subsection{Promises and certificates of convexity}
\label{sec:models-promises}

A function $f$ is \emph{convex on} a convex set $K$ if
$f(\lambda x+(1-\lambda)y)\le\lambda f(x)+(1-\lambda)f(y)$ for all $x,y\in K$
and $\lambda\in[0,1]$. It is \emph{globally convex} if it is convex on $\R^n$,
equivalently $\nabla^2f(x)\succeq0$ for all $x$. For $\mu>0$ it is
\emph{globally $\mu$-strongly convex} if $\nabla^2f(x)\succeq\mu I$ for all
$x\in\R^n$. A map $T:\R^n\to\R^n$ is \emph{globally $\mu$-strongly monotone}
if $(T(x)-T(y))^{\mathsf T}(x-y)\ge\mu\norm{x-y}^2$ for all $x,y$; for a
polynomial map this is equivalent to
$\frac12(J_T(x)+J_T(x)^{\mathsf T})\succeq\mu I$ for all $x$, where $J_T$ is
the Jacobian. A solution of the variational inequality defined by $T$ on a
closed convex set $K$ is a point $p\in K$ with
$T(p)^{\mathsf T}(y-p)\ge0$ for all $y\in K$; for $T=\nabla f$ these are the
minimizers of $f$ on $K$.

These properties enter in two ways. As a \emph{promise}, the property is
assumed and not checked; an algorithm for a promise problem may answer
arbitrarily on inputs that violate it. Recognizing convexity is NP-hard
already for quartics on $\R^n$~\cite{AhmadiOlshevskyParriloTsitsiklis2013}
and for cubics on a box~\cite{AhmadiHall2020}, so in general the promise
cannot be checked in polynomial time unless $\mathrm P=\NP$. As a
\emph{checked certificate}, the input contains data
that prove the property and can be verified in polynomial time. The input is
then an ordinary language, and a malformed certificate is mapped to a fixed
answer.

\begin{definition}[Hessian Gram matrices]
\label{def:models-hessian-gram}
Let $f\in\Q[x_1,\ldots,x_n]$ have degree at most four, and put
$w(x,v)=(v,\,x\otimes v)\in\R^{n+n^2}$. A \emph{Hessian Gram matrix} of $f$
is a matrix $A\in\mathrm{Sym}_{n+n^2}$ such that
\[
v^{\mathsf T}\nabla^2f(x)\,v=w(x,v)^{\mathsf T}A\,w(x,v)
\]
as polynomials in $(x,v)$. It is a \emph{full positive definite Hessian
Gram} if $A\succ0$. For higher degree the same definition is used with a
vector $w(x,v)$ of monomials that are linear in $v$ and include
$v_1,\ldots,v_n$ among them.
\end{definition}

A rational Hessian Gram matrix is checked in polynomial time by comparing
coefficients and by rational symmetric Gaussian elimination. Hessian Gram
matrices are not unique, since several entries contribute to the same
monomial of the biform.

\begin{lemma}[Determinant--trace margin]
\label{lem:models-det-trace}
Let $A\in\mathrm{Sym}_m$ be rational and positive definite, and put
\[
\rho(A)=\frac{\det A}{(\operatorname{tr}A)^{m-1}}.
\]
Then $A\succeq\rho(A)I$. The rational number $\rho(A)$ is computable in time
polynomial in the bit length of $A$ and has polynomial bit length.
\end{lemma}

\begin{proof}
Let $\lambda_1\le\cdots\le\lambda_m$ be the eigenvalues. Each is at most
$\operatorname{tr}A$, so
$\det A=\lambda_1\prod_{i\ge2}\lambda_i\le\lambda_1(\operatorname{tr}A)^{m-1}$.
The determinant of a rational matrix is computed in polynomial time with
polynomial bit length by fraction-free elimination, and the power
$(\operatorname{tr}A)^{m-1}$ has bit length at most $m$ times that of
$\operatorname{tr}A$.
\end{proof}

\begin{lemma}[Consequences of a Hessian Gram]
\label{lem:models-gram-curvature}\label{lem:models-hessian-gram}
Let $A$ be a Hessian Gram matrix of a quartic $f$ in $n$ variables.
\begin{enumerate}[label=(\alph*)]
\item If $A\succeq\mu I$ for some $\mu>0$, then $\nabla^2f(x)\succeq\mu I$ for
  all $x\in\R^n$. In particular, a full positive definite Hessian Gram
  certifies global strong convexity with the rational modulus $\rho(A)$ of
  Lemma~\ref{lem:models-det-trace}.
\item Let $T\in\Q^{n\times n}$ be invertible, $c\in\Q^n$, and
  $g(y)=f(Ty+c)$. Then $S^{\mathsf T}AS$ is a Hessian Gram matrix of $g$,
  where
  \[
  S=\begin{pmatrix}T&0\\ c\otimes T&T\otimes T\end{pmatrix}
  \]
  is invertible. Hence a full positive definite Hessian Gram is preserved by
  invertible rational affine changes of variables.
\item If $A\succ0$, then the homogeneous quartic part $f_4$ of $f$ satisfies
  $f_4(x)>0$ for all $x\ne0$.
\end{enumerate}
\end{lemma}

\begin{proof}
(a) Since $v$ is a subvector of $w=w(x,v)$,
$v^{\mathsf T}\nabla^2f(x)v=w^{\mathsf T}Aw\ge\mu\norm w^2\ge\mu\norm v^2$.
(b) We have $v^{\mathsf T}\nabla^2g(y)v=(Tv)^{\mathsf T}\nabla^2f(Ty+c)(Tv)$
and
$w(Ty+c,Tv)=(Tv,\,(T\otimes T)(y\otimes v)+(c\otimes T)v)=S\,w(y,v)$. The
matrix $S$ is block lower triangular with invertible diagonal blocks $T$ and
$T\otimes T$. (c) Let $Q$ be the block of $A$ indexed by $x\otimes v$.
Comparing the terms of degree two in $x$ in the defining identity gives
$v^{\mathsf T}\nabla^2f_4(x)v=(x\otimes v)^{\mathsf T}Q(x\otimes v)$. Setting
$v=x$ and using Euler's identity
$x^{\mathsf T}\nabla^2f_4(x)x=12f_4(x)$ gives
$12f_4(x)=(x\otimes x)^{\mathsf T}Q(x\otimes x)>0$ for $x\ne0$, because
$Q\succ0$.
\end{proof}

\begin{remark}[Hessian Gram and polynomial Gram matrices]
\label{rem:models-gram-kinds}
A full positive definite Hessian Gram is stronger than global strong
convexity. For example, $f(x,y)=x^2+y^2+y^4$ satisfies
$\nabla^2f\succeq2I$, but its quartic part $y^4$ vanishes at $(1,0)$, so by
Lemma~\ref{lem:models-hessian-gram}(c) it has no full positive definite Hessian
Gram. A sum $f(x)+g(y)$ in disjoint blocks of variables has none either: if
$v_i$ is a direction coordinate of the $y$-block, the monomial $x_j^2v_i^2$
has coefficient zero in the Hessian biform, and in $w^{\mathsf T}Aw$ it arises
only from the diagonal entry indexed by the coordinate $x_jv_i$ of $w$, which
must therefore vanish. A Hessian Gram certifies convexity, not
nonnegativity. A full positive definite Hessian Gram is compatible with
minimum zero at an irrational point, as for the quartic of
Section~\ref{sec:intro-optimizers}, whereas a positive definite polynomial
Gram matrix on the full monomial vector of
Definition~\ref{def:models-sos}, which contains $1$, implies strict
positivity at every real point. A positive definite Gram on a custom
polynomial vector can coexist with zeros when all entries of that vector
vanish there, as in Section~\ref{sec:fields-certificates}. The Hessian basis
$w(x,v)=(v,x\otimes v)$ instead contains every $v_i$, so it cannot vanish
when $v\ne0$; this is what makes its positive definiteness a global curvature
certificate.
\end{remark}

\begin{lemma}[Strong convexity]
\label{lem:models-strong-convexity}
Let $f$ be continuously differentiable and globally $\mu$-strongly convex,
and let $K\subseteq\R^n$ be nonempty, closed, and convex. Then $f$ attains its
minimum over $K$ at a unique point $p$, and for every $x\in K$
\[
f(x)-f(p)\ge\frac\mu2\norm{x-p}^2,
\qquad
\norm{x-p}\le\frac2\mu\norm{\nabla f(x)}.
\]
If $K=\R^n$, then also $\norm{x-p}\le\norm{\nabla f(x)}/\mu$.
\end{lemma}

\begin{proof}
For $x_0\in K$, $f(x)\ge f(x_0)+\nabla f(x_0)^{\mathsf T}(x-x_0)
+\frac\mu2\norm{x-x_0}^2$, so the sublevel sets of $f$ on $K$ are compact and
the minimum is attained; strict convexity makes the minimizer unique. First
order optimality gives $\nabla f(p)^{\mathsf T}(x-p)\ge0$, and Taylor's
formula with $\nabla^2f\succeq\mu I$ gives the first inequality. For the
second,
$0\ge f(p)-f(x)\ge\nabla f(x)^{\mathsf T}(p-x)+\frac\mu2\norm{p-x}^2
\ge-\norm{\nabla f(x)}\norm{p-x}+\frac\mu2\norm{p-x}^2$. If $K=\R^n$, then
$\nabla f(p)=0$, and
$\mu\norm{x-p}^2\le(\nabla f(x)-\nabla f(p))^{\mathsf T}(x-p)
\le\norm{\nabla f(x)}\norm{x-p}$.
\end{proof}

\begin{lemma}[Low degree]
\label{lem:models-low-degree}
\leavevmode
\begin{enumerate}[label=(\alph*)]
\item A globally convex polynomial of degree at most three has degree at most
  two. Consequently, if a globally convex polynomial of degree at most three
  attains its minimum on a rational polyhedron, then it has a rational
  minimizer whose bit length is polynomial in the input length, computable in
  polynomial time~\cite{KozlovTarasovKhachiyan1980}.
\item A globally monotone polynomial map $T:\R^n\to\R^n$ of degree at most
  two is affine. If it is strongly monotone and has rational coefficients, its
  unique zero is rational and is computed by rational linear algebra.
\end{enumerate}
\end{lemma}

\begin{proof}
(a) Write $f=f_0+f_1+f_2+f_3$ by homogeneous degree. For fixed $u$, the
function $x\mapsto u^{\mathsf T}\nabla^2f(x)u$ is affine and nonnegative on
$\R^n$, hence constant. Thus $u^{\mathsf T}\nabla^2f_3(x)u=0$ for all $x$ and
$u$, and Euler's identity gives $6f_3(x)=x^{\mathsf T}\nabla^2f_3(x)x=0$.
(b) Monotonicity is equivalent to $v^{\mathsf T}J_T(x)v\ge0$ for all $x$ and
$v$. For fixed $v$ this function of $x$ is affine and nonnegative, hence
constant. Let $U$ be the homogeneous quadratic part of $T$; then
$v^{\mathsf T}J_U(x)v=0$ for all $x,v$, that is,
$\partial_iU_j+\partial_jU_i=0$ for all $i,j$. Put
$a_{kij}=\partial_k\partial_iU_j$, a constant tensor symmetric in its first two
indices. Differentiating the identity gives $a_{kij}=-a_{kji}$. A tensor that
is symmetric in its first two indices and antisymmetric in its last two
vanishes:
$a_{kij}=-a_{kji}=-a_{jki}=a_{jik}=a_{ijk}=-a_{ikj}=-a_{kij}$. Hence all second
derivatives of $U$ vanish, so $U=0$ and $T(x)=Ax+c$. Strong monotonicity
gives $\frac12(A+A^{\mathsf T})\succ0$, so $A$ is invertible and
$p=-A^{-1}c$ is rational.
\end{proof}

Lemma~\ref{lem:models-low-degree} explains the degrees in this paper. For
globally convex objectives on rational polyhedra, degree four is the first
degree that permits a unique irrational optimizer. A convex cubic problem on
a rational polyhedron can have a unique irrational optimizer only if the
objective is convex on the feasible set but not on $\R^n$. Nonunique optimal
sets can contain irrational points already for quadratics; for example,
$f(x,y)=x^2$ is minimized at every point $(0,y)$.

\subsection{Output contracts}
\label{sec:models-outputs}

Let $f$ and $P$ be given, with $f_*$ finite and attained, and let
$S=\argmin_Pf$. The outputs studied in this paper are the following.

\begin{definition}[Values]
\label{def:models-value}
A \emph{value enclosure} of precision $q$ consists of a rational number
$\ell$ and a feasible rational point $x\in P$ with
$\ell\le f_*\le f(x)\le\ell+2^{-q}$. A \emph{value-gap point} is a feasible
rational $x$ with $f(x)-f_*\le2^{-q}$.
\end{definition}

\begin{definition}[Points]
\label{def:models-points}
A \emph{set approximation} of precision $q$ is a feasible rational $x$ with
$\dist(x,S)\le2^{-q}$. A \emph{fixed selector} consists of a point $p\in S$
determined by the instance, independently of $q$, and an algorithm that on
input $q$ returns a feasible rational $x_q$ with $\norm{x_q-p}\le2^{-q}$.
For convex problems the default selected point is the point of $S$ of minimum
Euclidean norm in the original coordinates. The sequence $(x_q)$ is a
Cauchy name of $p$; it need not give a short exact description of $p$.
\end{definition}

\begin{definition}[Exact predicates]
\label{def:models-predicates}
Let $h$ be an explicit rational polynomial (the observable). An \emph{exact
comparison} asks whether $h(p)\bowtie0$, for a relation
${\bowtie}\in\{<,\le,=,\ne,\ge,>\}$. Value thresholds ($h=f-r$), coordinate
thresholds ($h=x_j-r$), and constraint slacks ($h=b_i-B_ix$) are special
cases. The \emph{active set} of $p$ is the set of constraints with zero
slack. A \emph{Boolean predicate} is a polynomial-size Boolean combination of
exact comparisons. We call $<$ and $>$ strict, $\le$ and $\ge$ weak, and $=$
and $\ne$ equality relations.
\end{definition}

\begin{definition}[Exact representations]
\label{def:models-representations}
An exact representation of $p$ is one of the following.
\begin{enumerate}[label=(\alph*)]
\item An \emph{implicit} representation: a rational affine map $t\mapsto u+Mt$
  and an explicit polynomial system with a unique real solution $t_*$, such
  that $p=u+Mt_*$.
\item A \emph{circuit} representation of a rational point: a shared rational
  circuit (Definition~\ref{def:models-circuits}) whose value is $p$.
\item A \emph{root-circuit} representation, when available: a shared rational
  circuit with the explicitly defined real-root gates below. This format
  includes radical expressions but does not represent every algebraic point.
\item An \emph{algebraic} representation: for each coordinate, or for a
  primitive element of $\Q(p)$ together with the coordinates written as
  polynomials in it, the minimal polynomial over $\Q$ and an isolating
  interval. The minimal polynomial is given either as a \emph{dense list} of
  all its coefficients or as a \emph{sparse list} of its nonzero coefficients
  with their exponents.
\item An \emph{expanded rational} representation: reduced fractions for all
  coordinates.
\end{enumerate}
\end{definition}

These outputs form a hierarchy only in part. A fixed selector gives set
approximations and, by continuity, value-gap points; the converse fails,
because a set approximation may approach different optimizers at different
precisions. No approximation of fixed precision decides an exact
predicate, since $h(p)$ may be nonzero but smaller than any fixed bound. An
exact decision procedure need not produce a short representation, and a
short representation in one format need not be short in another: the number
$2^{1/d}$ has the two-term minimal polynomial $T^d-2$, while $1+2^{1/d}$ has
the minimal polynomial $(T-1)^d-2$, all of whose $d+1$ coefficients are
nonzero.

\subsection{Computational models}
\label{sec:models-computation}

An \emph{ordinary} algorithm is a deterministic algorithm in the bit model,
without an oracle. A \emph{Las Vegas} algorithm is randomized, returns a
correct answer whenever it halts, halts with probability one, and is
analyzed by its expected running time on each input. A parameterized
algorithm is \emph{fixed-parameter tractable} (FPT) if its running time is
at most $F(k)L^C$ with $F$ computable and $C$ an absolute constant; a bound
$L^{F(k)}$ is not of this form.

\begin{definition}[Arithmetic circuits]
\label{def:models-circuits}
An \emph{integer arithmetic circuit} is a finite directed acyclic graph
whose sources are labeled by the constants $0$ or $1$, and whose other nodes
are labeled by $+$, $-$, or $\times$ and have two ordered inputs; one node is
the output. Its size is the number of nodes, and nodes may feed several later
nodes. A circuit of size $s$ can compute an integer with about $2^s$ bits. A
\emph{rational circuit} may in addition use rational constants written in
binary and division by nonzero values. It is evaluated as a pair of integer
circuits $(N,M)$ with $M>0$, representing $N/M$: sums and products follow the
usual fraction rules, and division by a value $N_2/M_2\ne0$ uses
\[
\frac{N_1/M_1}{N_2/M_2}=\frac{N_1M_2N_2}{M_1N_2^2},
\]
which keeps the denominator positive without any sign test
(Lemma~\ref{lem:denominator-clearing}). A \emph{shared rational circuit} for
a vector or matrix is a single circuit with several outputs. Such circuits
have rational values. A \emph{root circuit} may also have a gate
$a\mapsto\sqrt[d]{a}$ with the positive integer $d$ written in binary: it
selects the unique real root when $d$ is odd and the nonnegative real root
when $d$ is even, in the latter case requiring $a\ge0$. A circuit with
algebraic inputs instead uses explicitly represented algebraic numbers as
sources. The representation of each such source is part of its size; a
rational circuit applied to an unspecified irrational source is not an exact
description.
\end{definition}

\begin{definition}[$\PosSLP$ and reductions]
\label{def:models-posslp}
$\PosSLP$ is the set of integer arithmetic circuits whose value is
positive~\cite{AllenderEtAl2009}. A problem \emph{reduces to one $\PosSLP$
instance} if a deterministic polynomial-time algorithm maps each instance $x$
to a circuit $C_x$ with $x$ a yes-instance exactly when $C_x\in\PosSLP$; for a
promise problem, the equivalence is required on instances that satisfy the
promise. The problem is \emph{$\PosSLP$-hard} if $\PosSLP$ reduces to it by
such a polynomial-time many-one reduction, and \emph{$\PosSLP$-complete} if
both hold. $\mathrm P^{\PosSLP}$ is the class of problems decided in
deterministic polynomial time by an oracle machine whose queries are circuit
descriptions. $\mathrm{UP}^{\PosSLP}$ is the class of problems decided by a
nondeterministic polynomial-time oracle machine with exactly one accepting
computation on each yes-instance and none on each no-instance, and
$\mathrm{coUP}^{\PosSLP}$ is the class of their complements. Las Vegas and FPT
algorithms with a $\PosSLP$ oracle are defined analogously.
\end{definition}

The sign of a rational circuit $(N,M)$ is the sign of $N$, so comparing a
rational circuit with zero is one $\PosSLP$ instance. Since $1-N>0$ exactly
when $N\le0$, strict and weak comparisons, and a problem and its complement,
reduce to each other. $\PosSLP$ lies in the counting
hierarchy~\cite[Theorem~1.4]{AllenderEtAl2009}, the class $\mathrm P^{\PosSLP}$
coincides with the class of Boolean problems decided in polynomial time by
constant-free real machines~\cite[Proposition~1.1]{AllenderEtAl2009}, and
Square Root Sum lies in
$\mathrm P^{\PosSLP}$~\cite[Section~1.4]{AllenderEtAl2009}.
Section~\ref{sec:upper} proves
that every problem in $\mathrm P^{\PosSLP}$ reduces to one $\PosSLP$ instance;
we therefore state deterministic oracle algorithms as single-instance
reductions. This conversion does not apply to $\mathrm{UP}^{\PosSLP}$, to Las
Vegas algorithms, or to algorithms whose running time is FPT rather than
polynomial.

These notions should be read with care. $\PosSLP$-hardness is hardness
relative to the sign problem of integer circuits. It is not NP-hardness, and
it does not exclude a polynomial-time algorithm unconditionally; whether
$\PosSLP$ is in P, or is NP-hard, is open~\cite{BuergisserJindal2024}.
Conversely, a reduction to one $\PosSLP$ instance is not an ordinary
algorithm: the circuit computes an integer whose binary expansion can have
exponentially many bits.

\begin{remark}[Exact conic feasibility]
\label{rem:models-socp}
Exact feasibility of second-order cone programs is $\PosSLP$-hard. This
follows from the basis equivalence of Tarasov and
Vyalyi~\cite[Theorem~3]{TarasovVyalyi2008} and standard conic duality, as
follows. An integer $N$ is positive exactly when $N\ge1$, so their basis
equivalence reduces $\PosSLP$ to deciding $A\ge B$ for the values $A,B$ of two
circuits that start from the constant $1$ and use only the operations $x+y$
and $x^2/2$; all gate values are positive. For one such
circuit, introduce a variable $x_i$ for every gate other than the input,
substitute $1$ for the input, and impose $x_i\ge x_j+x_k$ for an addition
gate and
\[
\begin{pmatrix}2x_i&x_j\\x_j&1\end{pmatrix}\succeq0,
\quad\text{that is,}\quad 2x_i\ge x_j^2,
\]
for a squaring gate. By induction in topological order, every feasible point
dominates the gate values, because addition is monotone and squaring is
monotone on positive numbers; the gate values themselves are feasible. Hence
the minimum of the output variable is the circuit value, and raising each
variable above its bound in topological order gives a strictly feasible
point. The cone is a product of half-lines and two-by-two positive
semidefinite cones, which is self-dual, so the dual program consists of
linear equations and membership in the same cone. Since the primal program
is strictly feasible and its value is finite, the conic duality theorem shows
that the dual attains the same value. Now require a dual feasible point for
the circuit of $A$, with objective value $d$, a primal feasible point for the
circuit of $B$, with objective value $e$, and the linear inequality $d\ge e$.
Weak duality gives $d\le A$ and $e\ge B$, so a feasible point forces $A\ge B$.
Conversely, if $A\ge B$, an optimal dual point with $d=A$ and the gate values
of the second circuit, with $e=B$, satisfy all constraints. The system is
therefore feasible exactly when $A\ge B$. Finally,
$\bigl(\begin{smallmatrix}a&b\\b&c\end{smallmatrix}\bigr)\succeq0$ if and only
if $\norm{(2b,a-c)}\le a+c$, so the system is a second-order cone program of
polynomial size. This consequence of~\cite{TarasovVyalyi2008} is recorded
for comparison only; the lower bounds of this paper concern a much smaller
class.
\end{remark}

\subsection{Certificates of nonnegativity}
\label{sec:models-certificates}

\begin{definition}[Gram matrices and sums of squares]
\label{def:models-sos}
Let $f$ have degree at most four, and let $z(x)$ be the vector of the
$\binom{n+2}{2}$ monomials of degree at most two, with the constant first.
A \emph{polynomial Gram matrix} of $f$ is a real symmetric matrix $Q$ with
$f=z^{\mathsf T}Qz$. Let $E\subseteq\R$ be a field. The polynomial $f$ is
a \emph{sum of squares (SOS) over $E$} if $f=\sum_ig_i^2$ with $g_i\in E[x]$,
and it has a \emph{positive semidefinite Gram matrix over $E$} if some
polynomial Gram matrix with entries in $E$ is positive semidefinite. A
positive definite polynomial Gram matrix is called \emph{interior}.
\end{definition}

An SOS over $E$ gives a positive semidefinite Gram matrix over $E$; the
converse holds for $E=\Q$ by Lemma~\ref{lem:models-rational-squares}, but not
for every real field~\cite{ChuaPlaumannSinnVinzant2017}, so field statements
for the two kinds of certificate need separate proofs. In an SOS
representation of a quartic, every $g_i$ has degree at most two, because the
highest-degree squares cannot cancel. If $f(p)=0$, every positive
semidefinite Gram matrix $Q$ satisfies $Qz(p)=0$, so $f$ has no interior
Gram matrix; every square factor then vanishes at $p$.

\begin{lemma}[Rational squares]
\label{lem:models-rational-squares}
Every positive rational $u/v$, with positive integers $u$ and $v$, is a sum
of at most $2\lceil\log_2(uv+1)\rceil$ squares of rationals, computable in time
polynomial in its bit length. Consequently, a rational positive semidefinite
Gram matrix of $f$ yields, in polynomial time, rational polynomials $g_i$ with
$f=\sum_ig_i^2$ and total bit length polynomial in that of the matrix.
\end{lemma}

\begin{proof}
Write $u/v=uv/v^2$ and expand $uv=\sum_j2^{e_j}$ in binary. Each term
$2^{2h}/v^2$ is the square $(2^h/v)^2$, and each term $2^{2h+1}/v^2$ is the sum
of two copies of that square. For the Gram matrix, rational symmetric
Gaussian elimination without square roots writes
$Q=\sum_id_i\ell_i\ell_i^{\mathsf T}$ with rationals $d_i\ge0$ and rational
vectors $\ell_i$; when a pivot vanishes, positive semidefiniteness forces its
row and column to vanish, and the row is skipped. Then
$f=\sum_id_i(\ell_i^{\mathsf T}z)^2$, and each $d_i$ is expanded as above.
The entries of the factors have polynomial bit length, as for fraction-free
elimination.
\end{proof}

Lemma~\ref{lem:models-rational-squares} avoids integer factorization and
four-square algorithms; it may use more rational squares than the rank of
$Q$, so it says nothing about the least number of rational squares.

\begin{definition}[Other certificate formats]
\label{def:models-formats}
\leavevmode
\begin{enumerate}[label=(\alph*)]
\item A polynomial $f$ is a \emph{sum of squares of rational functions with
  common denominator} $d\in\Q[x]\setminus\{0\}$ if $d^2f=\sum_ig_i^2$ with
  $g_i\in\Q[x]$. A \emph{multiplier certificate} is an identity
  $mf=\sum_ig_i^2$ with $m$ a polynomial that is positive on $\R^n$; a
  \emph{radial multiplier} is $m=(1+\norm{x}^2)^N$. The degree of a multiplier
  and the degree of a common denominator are different quantities.
\item For $f(x)+g(y)$ in disjoint blocks of variables, a \emph{block-separated
  certificate} over $E$ consists of a constant $c\in E$ and certificates over
  $E$ for $f+c$ and for $g-c$, each with its own monomial vector, including
  its own constant monomial.
\item The \emph{order-two moment relaxation} of a quartic $f$ minimizes
  $\sum_\alpha f_\alpha y_\alpha$ over sequences $(y_\alpha)_{\abs\alpha\le4}$
  with $y_0=1$ and moment matrix
    $M_2(y)=(y_{\alpha+\beta})_{\abs\alpha,\abs\beta\le2}\succeq0$. Its dual
  is the
  problem of maximizing $\gamma$ such that $f-\gamma$ has a positive
  semidefinite polynomial Gram matrix. When $f_*$ is attained in the dual,
  the positive semidefinite polynomial Gram matrices of $f-f_*$ are the
  \emph{optimal Gram matrices}. A nonzero positive semidefinite matrix $Z$
  with $\langle Z,Q\rangle=0$ for every real optimal Gram matrix $Q$ is an
  \emph{exposing matrix} for the set of optimal Gram matrices.
\end{enumerate}
\end{definition}

\subsection{Size measures}
\label{sec:models-size}

For an algebraic number $\alpha$, the degree is $[\Q(\alpha):\Q]$, and the
minimal polynomial is normalized to a primitive integer polynomial with
positive leading coefficient. For a point $p$, every coordinate satisfies
$[\Q(p_j):\Q]\le[\Q(p):\Q]$, so a lower bound on the degree of one
coordinate also bounds the degree of the joint field $\Q(p)$. The converse
fails: the joint field can have much larger degree than every individual
coordinate, so a lower bound on the joint degree need not force any single
coordinate to have comparable degree. The bit
length of a rational point is the sum of the bit lengths of its reduced
coordinates.

For input data in a number field, Appendix~\ref{app:further-arithmetic}
uses one common field $K=\Q(\alpha)$ with a specified real embedding. Its
defining polynomial, an isolating interval selecting $\alpha$, and the
coefficient lists expressing all data in the basis
$1,\alpha,\ldots,\alpha^{[K:\Q]-1}$ are part of the input. A witness may
require a larger field $K(p)$; its defining data and all coordinate lists
are charged to the output size. Separate short descriptions of individual
coordinates need not give a short dense description of their joint field.

Most of our lower-bound families have input length polynomial in the number
of variables $n$, or in a family index $k$, and output size exponential in
$n$ or $k$. Such bounds are superpolynomial in $L$, but they are not of the
form $2^{\Omega(L)}$, and we do not state them that way.

\subsection{Imported results}
\label{sec:models-imported}

Besides standard linear algebra, the proofs rely on the external results
listed in Table~\ref{tab:models-imported}. Each is stated with its exact
hypotheses where it is used.

{\small
\setlength{\tabcolsep}{4pt}
\setlength{\LTcapwidth}{\textwidth}
\begin{longtable}{@{}>{\raggedright\arraybackslash}p{0.52\textwidth}
  >{\raggedright\arraybackslash}p{0.28\textwidth}
  >{\raggedright\arraybackslash}p{0.14\textwidth}@{}}
\caption{External results used in the proofs, and where their precise
statements appear.}
\label{tab:models-imported}\\
\toprule
Result & Used for & Stated in\\
\midrule
\endfirsthead
\toprule
Result & Used for & Stated in\\
\midrule
\endhead
\bottomrule
\endfoot
Weak convex optimization by the ellipsoid method with separation oracles;
exact rational linear and convex quadratic
programming~\cite{GroetschelLovaszSchrijver1988,KozlovTarasovKhachiyan1980}
& Value approximation and starting points & Sec.~\ref{sec:points}\\
One-block real quantifier elimination with degree and coefficient-height
bounds~\cite[Theorem~2.27]{Basu2014} & Separation of nonzero values &
Secs.~\ref{sec:upper}, \ref{sec:constraints}\\
The rounded central-cut ellipsoid method~\cite{GroetschelLovaszSchrijver1988}
& Starting points for monotone maps & Sec.~\ref{sec:upper}\\
Linear programming with an operation count polynomial in the encoding of
the rational constraint matrix and independent of the objective and
right-hand-side encoding lengths~\cite{GroetschelLovaszSchrijver1988} & Normal-cone
verification & Sec.~\ref{sec:constraints}\\
Exact separable quadratic network flow with a strongly polynomial operation
count~\cite{Vegh2016}; exact constrained Newton
convergence~\cite{LeeSunSaunders2014} & Structured constrained problems &
Sec.~\ref{sec:constraints}\\
Violator-space sampling~\cite{GaertnerMatousekRuestSkovron2008}; integer
feasibility with separation queries at integer
points~\cite{AriHildebrand2026,HildebrandGoess2024,Basu2023Integers};
fixed-dimensional integer and mixed-integer linear
feasibility~\cite{HildebrandKoeppe2013,DelPia2023} & Rank and integer
parameters & Sec.~\ref{sec:constraints}\\
Certified isolation of real and complex roots of univariate
polynomials~\cite[Theorem~5]{MehlhornSagraloffWang2015} & Singleton
constructions & Sec.~\ref{sec:algebraic}\\
Bézout bounds for isolated solutions; Cayley--Bacharach and residual
intersection; classification of varieties of low degree; Brouwer
degree~\cite{Hartshorne1977,Fulton1998,EisenbudGreenHarris1996,
EklundJostPeterson2013,Harris1992,Milnor1965} & Degree bounds and number of
squares & Sec.~\ref{sec:algebraic}\\
Rational points of bounded height in open semialgebraic
sets~\cite[Theorem~4.1.2]{BasuPollackRoy1996} & Interior Gram matrices &
Sec.~\ref{sec:heights}\\
Characterization of the nonnegative rational ternary quartic forms that are
not sums of squares over $\Q$~\cite[Theorem~4.1]{Scheiderer2016} & Minimal
dimension of the field obstruction & Sec.~\ref{sec:fields}\\
Pure-state and order-unit criteria for sums of squares in archimedean
preorderings~\cite{BurgdorfScheidererSchweighofer2012} & Finite radial order &
Sec.~\ref{sec:fields-certificates}\\
Lattice basis reduction~\cite{LenstraLenstraLovasz1982}; two-block quantifier
elimination with separate dependence on coefficient
height~\cite{Renegar1992QE,BasuPollackRoy2006}; almost-everywhere
differentiability of convex conjugates and covering lemmas for maximal
functions~\cite{Rockafellar1970,EvansGariepy2015} & Core faces, margins,
exact fallback, and expected cell counts & Sec.~\ref{sec:recourse}\\
Multihomogeneous B\'ezout bounds; Hilbert irreducibility; algebraic-number
recognition~\cite{MorganSommese1987,Serre2008,KannanLenstraLovasz1988} &
Quadratic comparisons & App.~\ref{app:qc}\\
Effective quantifier elimination and {\L}ojasiewicz-type inequalities with
coefficient-height bounds~\cite{BasuMohammadNezhad2024} & Exact penalty
representations & App.~\ref{app:boundaries}\\
Sampling over quadratic maps with degree and height
bounds~\cite[Theorem~1.2]{GrigorievPasechnik2005} & Common-field witnesses
and finite infima & App.~\ref{app:further-arithmetic}\\
Integer-point bounds for convex first-order sets; rational mixed-integer
linear feasibility at fixed integer dimension~\cite{KhachiyanPorkolab2000,
Lenstra1983} & Algebraic cone feasibility &
App.~\ref{app:further-arithmetic}\\
\end{longtable}
}
